\documentclass[10pt,a4paper]{amsart}
\usepackage[margin=19mm]{geometry}
\usepackage{amsmath,amssymb,mathtools}
\usepackage[scr=rsfs,cal=euler]{mathalfa}
\usepackage{xfrac}
\usepackage[unicode,hidelinks,bookmarks=false]{hyperref}
\newcommand{\ie}{i.e.\ }
\newcommand{\cf}{cf.\ }
\newcommand{\resp}{respectively\ }
\newcommand{\Subsection}{\subsection}
\providecommand{\nobreakdash}{\nobreak\hbox{-}}
\setlength{\emergencystretch}{2em}
\multlinegap0pt
\usepackage{bm}
\usepackage{bbm}
\usepackage{dsfont}
\usepackage{xspace}
\usepackage{cancel}
\usepackage{mathtools}
\usepackage{amsthm}
\makeatletter
\@ifundefined{theorem}{\newtheorem{theorem}{Theorem}}{}
\makeatother
\def\h{\mathbb H}
\def\l{\mathbb L}
\def\r{\mathbb R}
\title[Extension of a problem of Euler in Lorentz-Minkowski plane]{Extension of a problem of Euler in Lorentz-Minkowski plane}
\author{Muhittin Evren Aydin, Rafael L\'opez}
\date{Source: arXiv:2508.17314v1 (CC BY 4.0)}
\begin{document}
\maketitle

\noindent\emph{Source and licence.} Extension of a problem of Euler in Lorentz-Minkowski plane. Version arXiv:2508.17314v1 (CC BY 4.0), distributed under the Creative Commons Attribution 4.0 International licence, as recorded in the arXiv licence metadata for this exact version and shown on the abstract page https://arxiv.org/abs/2508.17314v1. Full rights notice: licenses/arxiv-2508.17314-CC-BY-4.0.txt.\par\smallskip

\noindent\textit{Editorial scope.} Counted unit: the Theorem whose source label is \texttt{pr41} in the article (source0.tex lines 364--380, source label \texttt{pr41}), delivered together with the article's own complete original proof of it (source0.tex lines 382--417, one \texttt{proof} environment of 36 lines). No mathematics is written, repaired or re-derived: statement and proof are the article's own text line for line. Required context retained uncounted: k1 (lines 321--323). Every definition, display and label of those lines is retained unchanged. Omitted from this package and not redistributed: 1--320, 324--363, 381--381, 418--667. Nothing inside a retained excerpt is omitted or altered: every retained body is delivered exactly as the article's lines are written, so no omission receipt of any kind accompanies this package. Citations: the retained excerpts carry no citation command at all, so no bibliography entry of the article is transcribed and no reference is redistributed. Reference adaptation: none was needed and none was made; every reference inside the delivered unit resolves inside it, so no unresolved reference or citation is produced. Printed slips kept literal: no printed word, symbol, hypothesis or number of the counted statement or of its proof was altered. Layout and class: the article's own class is \texttt{amsart} with packages including amssymb, amsthm, graphicx, color, changes, url; the wrapper is the lane's pinned \texttt{amsart} editorial wrapper at 10pt on A4, which restates the theorem-like environments the retained text needs and adds the article's own macro definitions that the retained text needs (\texttt{\textbackslash h}, \texttt{\textbackslash l}, \texttt{\textbackslash r}), with extra wrapper packages (bm, bbm, dsfont, xspace, cancel, mathtools). The delivered numbering is the unit's own. Assets: no retained span contains a figure environment, a graphics-inclusion command, a PDF-page inclusion, a table or a TikZ picture; no image file is copied into this package, the package declares no asset dependency and no image file is redistributed with it. Licence: version arXiv:2508.17314v1 of this article is distributed under the Creative Commons Attribution 4.0 International licence, as recorded in the arXiv licence metadata for this exact version and shown on the abstract page https://arxiv.org/abs/2508.17314v1. A full rights notice is delivered at licenses/arxiv-2508.17314-CC-BY-4.0.txt. Characters: every retained excerpt is pure ASCII, so the delivered engine, XeLaTeX with its default Latin Modern text font, reports no missing character.
\begin{equation}\label{k1}
\rho  \rho ''+(\alpha -2) \rho '^2+(1-\alpha ) \rho ^2=0.
\end{equation}

\begin{theorem} \label{pr41}
Up to a linear isometry and a dilation of $\l^2$, the   $\alpha$-stationary spacelike curves contained in $\mathcal{C}^-$ are parametrized by $\gamma(s)=\rho(s)(\sinh(s),\cosh(s))$, where 
\begin{equation}\label{so}
\rho(s)=\left\{\begin{array}{ll}
 \cosh((\alpha-1) s)^\frac{1}{\alpha-1},&\alpha\not=1\\
 e^{c s},&\alpha=1, c^2<1.
 \end{array}
 \right.
\end{equation}
 Moreover, 
\begin{enumerate}
\item If $\alpha> 1$, then  the two branches of  $\gamma$ are asymptotic to a vertical translation of the lightlike cone $\mathcal{C}$.  
\item If $\alpha=1$, then $\gamma$ is a hyperbolic circle $\h^1(r)$ centered at the origin or the two branches of $\gamma$ are asymptotic to $\mathcal{C}$. 
\item If $0<\alpha<1$, then $\gamma$ meets tangentially $\mathcal{C}$ at two points.  
\item If $\alpha<0$, then $\gamma$ intersects orthogonally $\mathcal{C}$ at two points.
\end{enumerate}
\end{theorem}

\begin{proof}
Let $\gamma=\gamma(s)=(\gamma_1(s),\gamma_2(s))$ be an $\alpha$-stationary spacelike curve   included in $\mathcal{C}^-$. Let $\rho=\rho(s)$ be a smooth function such that  
$$\gamma(s)=\rho(s)(\sinh(s), \cosh(s)).$$
We know that $\rho$ satisfies \eqref{k1} with  $\rho^2-\rho'^2>0$. This equation   can be solved,  obtaining
$$\rho(s)= \left\{\begin{array}{ll}
c_2 \left(e^{(\alpha-1)s}+c_1 e^{-(\alpha-1)s}\right)^\frac{1}{\alpha-1},&\alpha\not=1\\
c_2e^{c_1 s},&\alpha=1, c_1^2<1.
\end{array}\right.
$$
where $c_1,c_2\in\r$, $c_2\not=0$. The spacelike condition on $\gamma$ implies  $c_1>0$ if $\alpha\not=1$ and $c_1^2<1$ if $\alpha=1$.  

We separate the case $\alpha\not=1$. Consider the rigid motion $R_t(x,y)= (\cosh(t) x+\sinh(t) y,\sinh(t)x+y\cosh(t))$, where $t=-\frac{\log(c_1)}{\alpha-1}$. Then 
$$R_t\cdot \gamma(s)=\gamma(s+t)=(2\sqrt{c_1}\cosh((\alpha-1)\theta))^\frac{1}{\alpha-1}(\sinh\theta,\cosh\theta),$$
where $\theta=s+t$. After a dilation, we obtain the parametrization given in \eqref{so}. Let 
$$\rho(s)=(\cosh((\alpha-1)s))^\frac{1}{\alpha-1}.$$
As a consequence, the curve $\gamma$ is symmetric about the $y$-axis.
\begin{enumerate}
\item Case $\alpha>1$. It is clear that $\lim_{s\to\infty}\gamma(s)=\infty$. In order to prove that it is asymptotic to a lightlike cone, we have
$$\lim_{s\to\pm\infty}\frac{\gamma_2'(s)}{\gamma_1'(s)}=\pm 1,$$
which gives the slope of the asymptote. Moreover, 
$$\lim_{s\to\infty}(\gamma_2(s)-\gamma_1(s))=2^\frac{1}{1-\alpha}.$$
This implies that the lines of equation $y=\pm  x+2^\frac{1}{1-\alpha}$ are asymptotic to $\gamma$. Thus $\gamma$ is asymptotic to $(0,2^{\frac{1}{1-\alpha}})+\mathcal{C}$ which is a vertical translation of $\mathcal{C}$.
\item Case $0<\alpha<1$. A computation gives 
$$\lim_{s\to\pm\infty}\gamma(s)=(\pm 2^\frac{\alpha}{1-\alpha},2^\frac{\alpha}{1-\alpha})$$
and
$$\lim_{s\to\pm\infty}\frac{\gamma_2'(s)}{\gamma_1'(s)}=\pm 1.$$
This proves that the intersection of $\gamma$ with  $\mathcal{C}$ is tangential.
\item Case $\alpha<0$. We have 
$$\lim_{s\to\pm\infty}\gamma(s)=(\pm 2^\frac{\alpha}{1-\alpha},2^\frac{\alpha}{1-\alpha}),$$
proving $\gamma$ tends to $\mathcal{C}$ at two points. Since we also have
$$\lim_{s\to\pm\infty}\frac{\gamma_2'(s)}{\gamma_1'(s)}=\pm 1,$$
the intersection (at the limit) between $\gamma$ and $\mathcal{C}$ is orthogonal.
\end{enumerate}
Suppose now $\alpha=1$. If $c_1=0$, then   $\gamma$ is the hyperbolic circle $\h^1(-c_2^2)$.  If $c_1\not=0$, then $\lim_{s\to\pm\infty}\frac{\gamma_2'(s)}{\gamma_1'(s)}=\pm 1$. However, now we have $\lim_{s\to\pm\infty}\gamma_2(s)-(\pm)\gamma_1(s)=0$. This proves that the two branches of $\gamma$ are asymptotic to $\mathcal{C}$.

\end{proof}

\end{document}
